\documentclass[12pt,a4paper]{article}

% Packages
\usepackage[utf8]{inputenc}
\usepackage[T1]{fontenc}
\usepackage{amsmath,amssymb,amsthm}
\usepackage{mathtools}
\usepackage{geometry}
\usepackage{graphicx}
\usepackage{booktabs}
\usepackage{array}
\usepackage{longtable}
\usepackage{hyperref}
\usepackage{cleveref}
\usepackage{algorithm}
\usepackage{algpseudocode}
\usepackage{listings}
\usepackage{xcolor}
\usepackage{natbib}
\usepackage{setspace}
\usepackage{fancyhdr}
\usepackage{titlesec}
\usepackage{enumitem}

% Page setup
\geometry{margin=1in}
\onehalfspacing

% Header/Footer
\pagestyle{fancy}
\fancyhf{}
\rhead{Research Proposal}
\lhead{Stochastic Optimization for Economic Forecasting}
\rfoot{Page \thepage}

% Colors
\definecolor{darkblue}{RGB}{31,78,121}
\definecolor{codegreen}{rgb}{0,0.6,0}
\definecolor{codegray}{rgb}{0.5,0.5,0.5}

% Hyperref setup
\hypersetup{
    colorlinks=true,
    linkcolor=darkblue,
    citecolor=darkblue,
    urlcolor=darkblue
}

% Code listing style
\lstset{
    basicstyle=\ttfamily\small,
    keywordstyle=\color{darkblue}\bfseries,
    commentstyle=\color{codegreen},
    stringstyle=\color{codegray},
    breaklines=true,
    frame=single
}

% Theorem environments
\newtheorem{definition}{Definition}
\newtheorem{proposition}{Proposition}
\newtheorem{theorem}{Theorem}

% Custom commands
\newcommand{\E}{\mathbb{E}}
\newcommand{\Var}{\mathrm{Var}}
\newcommand{\argmin}{\operatorname*{argmin}}
\newcommand{\R}{\mathbb{R}}

\title{
    \vspace{-1cm}
    \textbf{\Large Stochastic Optimization Methods for Robust Economic Forecasting Under Uncertainty} \\[0.5cm]
    \large Application to Serbian Composite Economic Index \\[1cm]
    \textsc{PhD/Master Research Proposal}
}

\author{
    \textbf{Candidate:} [Your Name] \\[0.3cm]
    \textbf{Proposed Mentor:} Prof. Dr. Dragan Vukmirovic \\
    Faculty of Organizational Sciences, University of Belgrade \\
    Department of Operational Research and Statistics \\[0.5cm]
    \texttt{dragan.vukmirovic@fon.bg.ac.rs}
}

\date{March 2026}

\begin{document}

\maketitle
\thispagestyle{empty}

\begin{abstract}
This research proposal presents a comprehensive framework for developing stochastic optimization methods applied to economic forecasting under uncertainty. Building on the seminal work on the Serbian Composite Index (SCI) developed by Prof.\ Vukmirovic and colleagues, this research extends the methodology by incorporating robust optimization techniques, Bayesian inference, and Monte Carlo simulation for uncertainty quantification in GDP nowcasting. The study utilizes real datasets from the Statistical Office of the Republic of Serbia (SORS), National Bank of Serbia (NBS), Eurostat, and FRED. We propose novel algorithms combining sample average approximation (SAA), stochastic gradient methods, and scenario-based optimization to construct robust composite indices that perform reliably across economic regimes.

\vspace{0.3cm}
\noindent\textbf{Keywords:} Stochastic Optimization, Economic Forecasting, Composite Index, GDP Nowcasting, Uncertainty Quantification, Bayesian Methods, Monte Carlo Simulation, Serbia
\end{abstract}

\newpage
\tableofcontents
\newpage

%==============================================================================
\section{Introduction and Motivation}
%==============================================================================

\subsection{Background}

Economic forecasting plays a critical role in policy-making, investment decisions, and business planning. Traditional econometric methods, while valuable, often fail to adequately capture the inherent uncertainty in economic systems, leading to overconfident predictions that can result in poor decisions during periods of economic stress.

The Serbian Composite Index (SCI), developed by Prof.\ Vukmirovic and colleagues at the Statistical Office of the Republic of Serbia \citep{vukmirovic2009sci, vukmirovic2013ntts}, represents a significant advancement in measuring short-term economic activity trends. The SCI provides monthly estimates of GDP dynamics, enabling policymakers to identify turning points in the business cycle more rapidly than traditional quarterly GDP releases.

However, current composite index methodologies typically assign fixed or equal weights to component indicators, which may not be optimal under different economic conditions. Moreover, they often lack rigorous uncertainty quantification, making it difficult to assess the reliability of forecasts.

\subsection{Research Gap}

While the I-distance methodology \citep{jovanovic2016idistance} provides a data-driven weighting scheme that overcomes biased weighting factors, significant opportunities remain to:

\begin{itemize}[leftmargin=*]
    \item Incorporate stochastic optimization to find weights robust to data uncertainty and model misspecification
    \item Develop adaptive weighting schemes responding to changing economic conditions
    \item Quantify forecast uncertainty through probabilistic methods
    \item Apply scenario-based optimization for stress testing
\end{itemize}

\subsection{Alignment with Mentor's Research}

This research directly builds upon Prof.\ Vukmirovic's extensive work in:
\begin{enumerate}[leftmargin=*]
    \item Serbian Composite Index for measuring economic activity (2009--present)
    \item GDP flash estimates and business cycle dating \citep{vukmirovic2015cycles}
    \item Multilevel I-distance methodology for composite indices \citep{jovanovic2016idistance}
    \item Official statistics and statistical system development
\end{enumerate}

%==============================================================================
\section{Research Objectives and Questions}
%==============================================================================

\subsection{Primary Research Objectives}

\begin{description}[leftmargin=!,labelwidth=1cm]
    \item[RO1] Develop a stochastic optimization framework for robust composite index construction accounting for data uncertainty, measurement error, and model misspecification.
    
    \item[RO2] Design adaptive weighting algorithms that adjust to different economic regimes (expansion, recession, crisis) using online learning and regime-switching models.
    
    \item[RO3] Implement Bayesian inference methods for uncertainty quantification in GDP nowcasting, providing credible intervals rather than point estimates.
    
    \item[RO4] Apply Monte Carlo simulation and scenario optimization for stress testing composite index performance under extreme conditions.
    
    \item[RO5] Validate proposed methods using historical Serbian economic data, including the 2008--2009 financial crisis and COVID-19 pandemic periods.
\end{description}

\subsection{Research Questions}

\begin{description}[leftmargin=!,labelwidth=1cm]
    \item[RQ1] How can stochastic optimization improve the robustness of composite economic indices compared to fixed-weight approaches?
    
    \item[RQ2] What is the optimal balance between forecast accuracy and uncertainty quantification in GDP nowcasting?
    
    \item[RQ3] How do adaptive weighting schemes perform during economic regime transitions?
    
    \item[RQ4] Can Bayesian methods provide reliable uncertainty bounds for composite index values?
    
    \item[RQ5] What computational trade-offs exist between algorithm complexity and real-time forecasting requirements?
\end{description}

%==============================================================================
\section{Literature Review}
%==============================================================================

\subsection{Composite Indices in Economic Measurement}

Composite indicators have become essential tools for measuring complex, multidimensional phenomena. The OECD Handbook on Constructing Composite Indicators \citep{nardo2005handbook} provides foundational methodology. The I-distance method, applied to the Networked Readiness Index \citep{jovanovic2016idistance}, addresses equal-weight limitations through discriminant analysis-based weighting.

\subsection{Stochastic Optimization in Economics}

Stochastic optimization has seen extensive application in finance but limited adoption in macroeconomic forecasting. Key methodologies include:

\begin{definition}[Sample Average Approximation]
Given a stochastic program $\min_x \E_\xi[f(x,\xi)]$, SAA replaces the expectation with a sample mean:
\begin{equation}
    \min_x \frac{1}{N}\sum_{n=1}^{N} f(x, \xi^n)
\end{equation}
where $\{\xi^1, \ldots, \xi^N\}$ are i.i.d.\ samples from the distribution of $\xi$.
\end{definition}

\begin{definition}[Robust Optimization]
Seeks solutions performing well across all scenarios in an uncertainty set $\mathcal{U}$:
\begin{equation}
    \min_x \max_{\xi \in \mathcal{U}} f(x, \xi)
\end{equation}
\end{definition}

\subsection{GDP Nowcasting and Business Cycles}

GDP nowcasting---estimating current-quarter GDP before official release---has gained prominence using dynamic factor models and machine learning. The Serbian Composite Index work \citep{vukmirovic2009sci, vukmirovic2010flash, vukmirovic2013ntts} demonstrated monthly GDP tracking feasibility in a transition economy.

%==============================================================================
\section{Datasets and Data Sources}
%==============================================================================

This research utilizes exclusively real, publicly available datasets. No synthetic data will be used.

\subsection{Primary Data Sources}

\begin{table}[htbp]
\centering
\caption{Overview of Data Sources}
\label{tab:datasources}
\begin{tabular}{@{}llll@{}}
\toprule
\textbf{Source} & \textbf{Key Variables} & \textbf{Frequency} & \textbf{Coverage} \\
\midrule
SORS & GDP, IPI, Retail Trade, CPI & Monthly/Quarterly & 1995--present \\
NBS & Interest rates, FX, Money supply & Daily/Monthly & 1999--present \\
FRED & Serbia GDP, Inflation & Quarterly/Annual & 1995--present \\
Eurostat & EU harmonized indicators & Monthly/Quarterly & 2006--present \\
World Bank & WDI, Structural indicators & Annual & 1990--present \\
\bottomrule
\end{tabular}
\end{table}

\subsection{Statistical Office of the Republic of Serbia (SORS)}

The SORS database (\url{https://data.stat.gov.rs}) provides core national statistics:

\begin{itemize}[leftmargin=*]
    \item \textbf{Quarterly GDP}: Real GDP at constant prices, chain-linked volumes, 1995--present
    \item \textbf{Industrial Production Index (IPI)}: Monthly, seasonally adjusted, base 2021=100
    \item \textbf{Retail Trade Turnover}: Monthly, current and constant prices
    \item \textbf{Consumer Price Index (CPI)}: Monthly, HICP-compatible methodology
    \item \textbf{Labour Force Survey}: Employment, unemployment rates, quarterly
\end{itemize}

\subsection{National Bank of Serbia (NBS)}

The NBS Statistical Bulletin (\url{https://nbs.rs/en/statistika/}) provides:

\begin{itemize}[leftmargin=*]
    \item \textbf{Key Policy Rate}: Currently 5.75\% (February 2026), daily series from 2007
    \item \textbf{Exchange Rates}: Daily RSD/EUR, RSD/USD
    \item \textbf{Money Supply}: M0, M1, M2, M3 aggregates, monthly
    \item \textbf{Inflation Expectations Survey}: Household and corporate expectations
\end{itemize}

\subsection{FRED (Federal Reserve Economic Data)}

FRED (\url{https://fred.stlouisfed.org}) hosts Eurostat-sourced Serbian data:

\begin{itemize}[leftmargin=*]
    \item \texttt{CLVMNACNSAB1GQRS}: Serbia Real GDP, quarterly, chain-linked volumes
    \item \texttt{FPCPITOTLZGSRB}: Serbia Inflation, annual CPI, 1995--2024
    \item EU economic indicators for external shock modeling
\end{itemize}

\subsection{Eurostat and World Bank}

\begin{itemize}[leftmargin=*]
    \item \textbf{Eurostat}: ESA 2010 methodology GDP, business/consumer surveys, trade statistics
    \item \textbf{World Bank}: World Development Indicators for structural analysis
\end{itemize}

%==============================================================================
\section{Proposed Methodology}
%==============================================================================

\subsection{Overall Framework}

The methodology integrates three complementary approaches:

\begin{table}[htbp]
\centering
\caption{Methodological Components}
\label{tab:methodology}
\begin{tabular}{@{}lll@{}}
\toprule
\textbf{Component} & \textbf{Methods} & \textbf{Purpose} \\
\midrule
Stochastic Optimization & SAA, Robust Opt., SGD & Weight determination \\
Bayesian Inference & MCMC, State-space & Uncertainty quantification \\
Monte Carlo Simulation & Scenarios, Sensitivity & Stress testing \\
\bottomrule
\end{tabular}
\end{table}

\subsection{Stochastic Optimization for Weight Determination}

\subsubsection{Problem Formulation}

Let $\mathbf{X} = \{X_1, X_2, \ldots, X_K\}$ denote $K$ component indicators and $\mathbf{w} = \{w_1, w_2, \ldots, w_K\}$ be the weight vector. The composite index is:

\begin{equation}
    C(\mathbf{w}, \mathbf{X}) = \sum_{k=1}^{K} w_k \cdot X_k, \quad \text{subject to: } \sum_{k=1}^{K} w_k = 1, \; w_k \geq 0
\end{equation}

The stochastic optimization problem seeks weights minimizing expected forecast error under uncertainty:

\begin{equation}
    \min_{\mathbf{w} \in \mathcal{W}} \E_{\boldsymbol{\xi}}\left[ L\left(C(\mathbf{w}, \mathbf{X} + \boldsymbol{\xi}), \text{GDP}\right) \right]
    \label{eq:stochastic_problem}
\end{equation}

where $\boldsymbol{\xi}$ represents random perturbations capturing measurement error, revision uncertainty, and model misspecification. $L(\cdot)$ is a loss function (e.g., MSE, quantile loss).

\subsubsection{Sample Average Approximation (SAA)}

We approximate the expected loss using $N$ Monte Carlo samples:

\begin{equation}
    \min_{\mathbf{w} \in \mathcal{W}} \frac{1}{N} \sum_{n=1}^{N} L\left(C(\mathbf{w}, \mathbf{X} + \boldsymbol{\xi}^{(n)}), \text{GDP}\right)
    \label{eq:saa}
\end{equation}

\begin{proposition}[SAA Convergence]
Under standard regularity conditions, as $N \to \infty$, the SAA optimal value converges almost surely to the true optimal value of \eqref{eq:stochastic_problem}.
\end{proposition}

\subsubsection{Robust Optimization Extension}

For crisis-period robustness, we incorporate an uncertainty set $\mathcal{U}$:

\begin{equation}
    \min_{\mathbf{w} \in \mathcal{W}} \max_{\boldsymbol{\xi} \in \mathcal{U}} L\left(C(\mathbf{w}, \mathbf{X} + \boldsymbol{\xi}), \text{GDP}\right)
    \label{eq:robust}
\end{equation}

The uncertainty set can be:
\begin{itemize}
    \item \textbf{Ellipsoidal}: $\mathcal{U} = \{\boldsymbol{\xi} : \boldsymbol{\xi}^\top \Sigma^{-1} \boldsymbol{\xi} \leq \Gamma^2\}$ based on historical covariance
    \item \textbf{Polyhedral}: $\mathcal{U} = \{\boldsymbol{\xi} : A\boldsymbol{\xi} \leq \mathbf{b}\}$ based on revision ranges
\end{itemize}

\subsubsection{Adaptive Stochastic Gradient Descent}

For online learning with streaming data, we employ adaptive SGD:

\begin{algorithm}[htbp]
\caption{Adaptive Composite Index Optimization}
\label{alg:adaptive_sgd}
\begin{algorithmic}[1]
\Require Initial weights $\mathbf{w}_0$, learning rate $\eta$, decay $\gamma$
\For{$t = 1, 2, \ldots, T$}
    \State Observe new indicator data $\mathbf{X}_t$
    \State Sample perturbation $\boldsymbol{\xi}_t \sim \mathcal{N}(\mathbf{0}, \hat{\Sigma}_t)$
    \State Compute gradient $\mathbf{g}_t = \nabla_{\mathbf{w}} L(C(\mathbf{w}_{t-1}, \mathbf{X}_t + \boldsymbol{\xi}_t), \text{GDP}_t)$
    \State Update: $\mathbf{w}_t = \Pi_{\mathcal{W}}\left(\mathbf{w}_{t-1} - \eta_t \mathbf{g}_t\right)$
    \State Decay: $\eta_{t+1} = \eta_t \cdot \gamma$
\EndFor
\State \Return $\mathbf{w}_T$
\end{algorithmic}
\end{algorithm}

\subsection{Bayesian Inference for Uncertainty Quantification}

\subsubsection{Hierarchical Bayesian Model}

We model the composite index--GDP relationship:

\begin{align}
    \text{GDP}_t &\sim \mathcal{N}(\alpha + \beta \cdot C_t, \sigma^2) \label{eq:likelihood} \\
    \alpha, \beta &\sim \mathcal{N}(0, \tau^2) \label{eq:prior_coef} \\
    \sigma &\sim \text{Half-Cauchy}(0, 1) \label{eq:prior_sigma}
\end{align}

Posterior distributions computed via MCMC provide full uncertainty quantification.

\subsubsection{Time-Varying Parameter Model}

To capture structural changes, we extend to state-space:

\begin{align}
    \text{GDP}_t &= \alpha_t + \beta_t \cdot C_t + \epsilon_t, \quad \epsilon_t \sim \mathcal{N}(0, \sigma^2) \\
    \beta_t &= \beta_{t-1} + \eta_t, \quad \eta_t \sim \mathcal{N}(0, Q)
\end{align}

This allows the relationship to evolve, capturing regime changes.

\subsubsection{Posterior Predictive Distribution}

For a new observation $C^*$, the predictive distribution is:

\begin{equation}
    p(\text{GDP}^* | C^*, \mathcal{D}) = \int p(\text{GDP}^* | C^*, \theta) \, p(\theta | \mathcal{D}) \, d\theta
\end{equation}

where $\mathcal{D}$ is observed data and $\theta = \{\alpha, \beta, \sigma\}$.

\subsection{Monte Carlo Simulation for Stress Testing}

\subsubsection{Scenario Generation}

We generate $S$ scenarios from estimated joint distribution:

\begin{equation}
    \mathbf{X}^{(s)} \sim \mathcal{N}(\hat{\boldsymbol{\mu}}, \hat{\Sigma}), \quad s = 1, \ldots, S
\end{equation}

For tail risk analysis, we condition on extreme quantiles:

\begin{equation}
    \mathbf{X}^{(s)} | \left(\mathbf{X}^{(s)} \in \mathcal{Q}_\alpha\right)
\end{equation}

where $\mathcal{Q}_\alpha$ is the $\alpha$-quantile region.

\subsubsection{Sensitivity Analysis}

Sobol indices quantify which indicators most affect uncertainty:

\begin{equation}
    S_k = \frac{\Var_{X_k}\left(\E_{\mathbf{X}_{-k}}[C | X_k]\right)}{\Var(C)}
\end{equation}

where $\mathbf{X}_{-k}$ denotes all indicators except $X_k$.

%==============================================================================
\section{Implementation Plan}
%==============================================================================

\subsection{Software Stack}

\begin{table}[htbp]
\centering
\caption{Software and Libraries}
\label{tab:software}
\begin{tabular}{@{}ll@{}}
\toprule
\textbf{Component} & \textbf{Tools} \\
\midrule
Core & Python 3.11+, NumPy, SciPy, Pandas \\
Optimization & CVXPY, SciPy.optimize, OR-Tools \\
Bayesian & PyMC, ArviZ, Stan \\
Time Series & statsmodels, arch, darts \\
Data Access & fredapi, eurostat (R/rpy2), requests \\
\bottomrule
\end{tabular}
\end{table}

\subsection{Research Timeline}

\begin{table}[htbp]
\centering
\caption{30-Month Research Timeline}
\label{tab:timeline}
\begin{tabular}{@{}cll@{}}
\toprule
\textbf{Phase} & \textbf{Activities} & \textbf{Months} \\
\midrule
1 & Literature review, data collection, EDA & 1--4 \\
2 & Stochastic optimization algorithm development & 5--10 \\
3 & Bayesian model implementation & 11--16 \\
4 & Monte Carlo stress testing, validation & 17--22 \\
5 & Software release, papers, thesis defense & 23--30 \\
\bottomrule
\end{tabular}
\end{table}

%==============================================================================
\section{Expected Contributions}
%==============================================================================

\subsection{Theoretical Contributions}

\begin{enumerate}[leftmargin=*]
    \item Novel integration of stochastic optimization with composite indicator methodology
    \item Formal uncertainty quantification framework for GDP nowcasting in transition economies
    \item Adaptive weighting algorithms with convergence guarantees under regime changes
\end{enumerate}

\subsection{Practical Contributions}

\begin{enumerate}[leftmargin=*]
    \item Open-source Python package for stochastic composite index construction
    \item Updated Serbian Composite Index with uncertainty bounds
    \item Stress testing toolkit for statistical offices and central banks
    \item Methodology transferable to Western Balkan economies
\end{enumerate}

\subsection{Target Publications}

\begin{itemize}[leftmargin=*]
    \item European Journal of Operational Research (methodology)
    \item Journal of Applied Econometrics / International Journal of Forecasting (application)
    \item Social Indicators Research / Information Development (composite indices)
    \item NTTS, IAOS conferences
\end{itemize}

%==============================================================================
\section{Conclusion}
%==============================================================================

This research proposal outlines a comprehensive framework for applying stochastic optimization to economic forecasting, directly building on Prof.\ Vukmirovic's foundational work on the Serbian Composite Index. By integrating robust optimization, Bayesian inference, and Monte Carlo methods, we aim to produce forecasts that are not only accurate but also honest about their uncertainty---a critical requirement for sound economic policy-making.

%==============================================================================
% References
%==============================================================================

\bibliographystyle{apalike}
\begin{thebibliography}{99}

\bibitem[Ben-Tal et al., 2009]{bental2009robust}
Ben-Tal, A., El Ghaoui, L., \& Nemirovski, A. (2009).
\newblock {\em Robust Optimization}.
\newblock Princeton University Press.

\bibitem[Jovanovic Milenkovic et al., 2016]{jovanovic2016idistance}
Jovanovic Milenkovic, M., Brajovic, B., Milenkovic, D., Vukmirovic, D., \& Jeremic, V. (2016).
\newblock Beyond the equal-weight framework of the Networked Readiness Index: A multilevel I-distance methodology.
\newblock {\em Information Development}, 32(4), 1120--1136.

\bibitem[Nardo et al., 2005]{nardo2005handbook}
Nardo, M., Saisana, M., Saltelli, A., \& Tarantola, S. (2005).
\newblock Handbook on Constructing Composite Indicators.
\newblock {\em OECD Statistics Working Papers}.

\bibitem[Rahimian \& Mehrotra, 2019]{rahimian2019dro}
Rahimian, H., \& Mehrotra, S. (2019).
\newblock Distributionally Robust Optimization: A Review.
\newblock {\em arXiv preprint arXiv:1908.05659}.

\bibitem[Shapiro et al., 2021]{shapiro2021lectures}
Shapiro, A., Dentcheva, D., \& Ruszczynski, A. (2021).
\newblock {\em Lectures on Stochastic Programming: Modeling and Theory}.
\newblock SIAM.

\bibitem[Vukmirovic et al., 2009]{vukmirovic2009sci}
Vukmirovic, D., Milojic, A., Ciric, R., \& Vukmirovic, J. (2009).
\newblock Serbian Composite Index as an instrument for measuring short-term trends of economic activity.
\newblock {\em Statistics Journal}, National Statistical Institute Bulgaria.

\bibitem[Vukmirovic \& Ciric, 2010]{vukmirovic2010flash}
Vukmirovic, D., \& Ciric, R. (2010).
\newblock Using the Serbian Composite Index for Flash Estimates of GDP.
\newblock {\em 6th Colloquium on Modern Tools for Business Cycle Analysis}, Eurostat.

\bibitem[Vukmirovic et al., 2013]{vukmirovic2013ntts}
Vukmirovic, D., et al. (2013).
\newblock Serbian Composite Index -- Indicator of Monthly GDP.
\newblock {\em NTTS 2013}, Eurostat.

\bibitem[Vukmirovic et al., 2015]{vukmirovic2015cycles}
Vukmirovic, D., Ciric, R., Smolcic, M., Jelic, S., \& Karamarkovic, S. (2015).
\newblock Dating the Serbian Business Cycles.
\newblock {\em Journal of Economics, Business and Management}, 3(6).

\end{thebibliography}

\end{document}
